\documentclass{EEIMVR}
\usepackage[utf8]{inputenc}
\usepackage[T1]{fontenc}
\usepackage{geometry}
\usepackage{graphicx}
\usepackage{float}
\usepackage{hyperref}
\usepackage[english,indonesian]{babel}

\title{Optimasi Pemrosesan Template LaTeX Menggunakan Pendekatan Berbasis AI}
\author{Academic Support Agent}
\date{\today}

\begin{document}
\maketitle

\begin{abstract}

\end{abstract}

\chapter{PENDAHULUAN}


\chapter{TINJAUAN PUSTAKA}


% Original content (commented out):
% \documentclass[a4paper,12pt,oneside,onecolumn,final]{EEIMVR}
% ---
% Pacotes fundamentais 
% ---

\usepackage{lipsum}
% ---
% Pacote auxiliar para as normas da UFF
% ---
\usepackage[frame=no,font=sans]{EEIMVRformat} %font=sans (arial), font=times, font=default ou font=lmodern

% ---
% Pacote de citacoes
% ---
\usepackage[alf, abnt-repeated-author-omit=no, abnt-emphasize = bf, abnt-full-initials=no, abnt-thesis-year=final]{abntex2cite}

% **********************************************************
% **********************************************************
% Informações de autoria e institucionais
% **********************************************************
% **********************************************************

%-----------------------------------------------------------
% Grau pretendido (Doutor, Mestre, Bacharel, Licenciado) e Curso
%-----------------------------------------------------------


\grau{Graduação}
% \grau{Mestre} 
%\grau{Doutor}


\curso{Engenharia Mecânica}


\areadeconcentracao{} %Projeto de Graduação
% \areadeconcentracao{Fenômenos de Transporte} %Pós Graduação
%\areadeconcentracao{Mecânica dos Sólidos} %Pós Graduação

%-----------------------------------------------------------
% Informações da instituição
%-----------------------------------------------------------
\instituicao{Universidade Federal Fluminense}  %Universidade
            {} % Polo Universitário
            {Escola de Engenharia Industrial Metalúrgica de Volta Redonda} %Unidade
            {Departamento de \UFFcurso}

%-----------------------------------------------------------
% Informações da autoria do documento
%-----------------------------------------------------------

\autor{Nome e Sobrenome do Autor do}
      {Trabalho}
    %   {iniciais. do. nome.} % iniciais do nome

\titulo{Título do trabalho acadêmico em português} %Título do trabalho acadêmico

\orientador{Prof.} %cargo, ex.: Prof., Profa., Eng. , etc...
		   {Jorge}{Alberto Rodriguez Durán, D.Sc.} %Nome sobrenome com a titulação ao final. Exemplo: D.Sc., Ph.D., M.Sc., B.Sc., etc.
           {UFF -- Universidade Federal Fluminense -- VEM} % Instituição e Departamento ou PPG


%Opcional, Comente as linhas de coorientador caso não tenha
\coorientador{cargo}  %cargo, ex.: Prof., Profa., Eng., etc
           {nome}{sobrenome, titulação}  %Nome sobrenome com a titulação ao final. Exemplo: D.Sc., Ph.D., M.Sc., B.Sc., etc.
           {unidade -- instituição} % Departamento ou PPG e Instituição 

%-----------------------------------------------------------
% Informações adicionais (local, data e paginas)
%-----------------------------------------------------------

\local{Volta Redonda, RJ} 
\data{DD}{MMMMMM}{AAAA} % Exemplo: \data{21}{Março}{2016}

% ***********************************************************
% ***********************************************************
% Início do documento
% ***********************************************************
% ***********************************************************

\begin{document}
% ----------------------------------------------------------
%% ELEMENTOS PRE-TEXTUAIS
% ----------------------------------------------------------
\frontmatter
% ----------------------------------------------------------
% Capa e a folha de rosto
% ----------------------------------------------------------
\capa % use \capa* para ir para o seguundo modelo de capa
\folhaderosto
% ----------------------------------------------------------
% Inserir a ficha catalográfica
% ----------------------------------------------------------
% ---
% A biblioteca deverá providenciar a ficha catalográfica. Salve a ficha no
% formato PDF. Use o nome do arquivo PDF como argumento do comando. 
% Exemplo: ficha catalográfica é o arquivo 'Ficha.pdf' na pasta "B.PreTextual"
%     \fichacatalografica{B.PreTextual/Ficha.pdf}
%
% Enquanto não possuir a ficha catalográfica, use o comando sem argumentos.
% ---
% \includepdf[noautoscale, pagecommand={}]{B.PreTextual/Ficha.pdf} % Substituir o arquivo ficha.pdf pela sua ficha catalográfica

\include{B.PreTextual/ficha}



% ----------------------------------------------------------
% Folha de aprovação
% ----------------------------------------------------------

%Após obter a assinatura dos membros da banca, comente as linhas a baixos e insira o pdf com as assinaturas na pasta "B.PreTextual"


% Ajustar o espaço entre as assinaturas abaixo
\newcommand{\spc}{0.25cm}

\begin{folhadeaprovacao}
	\vspace{\spc} % Espaço entre assinaturas
	\assinatura{Engenheiro João Carlos de Nascimento, M.Sc.}{instituição -- Alguma empresa da vida} %Exemplo
% 	\vspace{\spc} % Espaço entre assinaturas
% 	\assinatura{Eng. João José, M.Sc.}{Instituição}
% 	\vspace{\spc} % Espaço entre assinaturas
% 	\assinatura{Eng. Marcos João, B.Sc.}{Instituição}
% 	\vspace{\spc} % Espaço entre assinaturas
% 	\assinatura{quarto membro titular da banca}{Instituição}
\end{folhadeaprovacao}


% Após colocar o pdf com as assinaturas na pasta "B.PreTextual", comente todo o ambiente "folhadeaprovacao" acima, descomente a linha abaixo e insira o nome correto do arquivo pdf:
%\includepdf[pages=1]{B.PreTextual/ficha.pdf} %exemplo


% ----------------------------------------------------------
% Dedicatória
\include{B.PreTextual/01.Dedication}

% ----------------------------------------------------------
% Agradecimentos
\include{B.PreTextual/02.Acknowledgment}

% ----------------------------------------------------------
% Epigrafe (opcional)
\include{B.PreTextual/03.Epigraph}

% ----------------------------------------------------------
%% RESUMO
\include{B.PreTextual/04.Abstract_pt}

% ----------------------------------------------------------
% Abstract
\begin{otherlanguage}{english}
\include{B.PreTextual/05.Abstract_en}
\end{otherlanguage}



% ----------------------------------------------------------
% Listas de ilustrações e tabelas
% ----------------------------------------------------------
\listoffigures % list of figures
\listoftables % list of tables
% ----------------------------------------------------------
% Outras listas
% ----------------------------------------------------------
%\listadealgoritmos


% ----------------------------------------------------------
% Lista de abreviaturas e siglas
\include{B.PreTextual/06.AbbreviationList}

% ----------------------------------------------------------
% Lista de simbolos
\include{B.PreTextual/07.SymbolList}

% ----------------------------------------------------------
% Sumario
\tableofcontents % table of contents

% ----------------------------------------------------------
%% ELEMENTOS TEXTUAIS
% ----------------------------------------------------------
\mainmatter



\include{C.Chapters/01.Introduction}
\include{C.Chapters/02.Theoretical_Background}
\include{C.Chapters/90.Results}
\include{C.Chapters/99.Conclusion}




% ----------------------------------------------------------
%% ELEMENTOS POS-TEXTUAIS
% ----------------------------------------------------------
\backmatter
%=====================================================================
% Referencias via BibTeX
%=====================================================================

\bibliography{A.Bibliography/MyBibliography}
%=====================================================================




%=====================================================================
%% Glossário
%=====================================================================
\include{D.PostTextual/01.Glossary}




% ----------------------------------------------------------
% Apêndices (opcionais)
% ----------------------------------------------------------
% ---
% Inicia os apêndices
% ---
\appendix

\include{D.PostTextual/02.Appendix-A}

\include{D.PostTextual/02.Appendix-B}






% ----------------------------------------------------------
% Anexos (opcionais)
% ----------------------------------------------------------
% ---
% Inicia os anexos
% ---
\annex

\include{D.PostTextual/03.Annex-A}

\include{D.PostTextual/03.Annex-B}

%=====================================================================
\end{document}


\end{document}